% psmc-rs as a GigaByte Technical Release, in the class of GigaByte's Overleaf template ("Template for
% GigaByte Journal Technical Report Submissions", oup-contemporary; its files are in gigabyte-template/).
% `make gigabyte` collects everything in build/gigabyte/ and compiles it there; the same files compile on
% Overleaf. Without oup-contemporary.cls the file still compiles, as a plain preview. Every result quoted is a macro of numbers.tex, generated by
% make_numbers.py from the committed results, as in the article (psmc-rs.tex).
\IfFileExists{oup-contemporary.cls}{%
  \documentclass[a4paper,num-refs,gigabyte]{oup-contemporary}%
}{%
  \documentclass[11pt,a4paper]{article}%
  \def\gigabytepreview{}%
}

\ifdefined\gigabytepreview
  % A plain stand-in for oup-contemporary, so that the file also compiles without it.
  \usepackage[a4paper,margin=2cm]{geometry}
  \usepackage[numbers,sort&compress]{natbib}
  \bibliographystyle{unsrtnat}
  \usepackage{authblk}
  \usepackage{amsmath}
  \usepackage[hidelinks]{hyperref}
  \usepackage{url}
  \newcommand{\journal}[1]{}
  \newcommand{\papercat}[1]{}
  \newcommand{\runningauthor}[1]{}
  \newcommand{\jvolume}[1]{}
  \newcommand{\jnumber}[1]{}
  \newcommand{\jyear}[1]{}
  \newcommand{\authfn}[1]{*}
  \newcommand{\authnote}[1]{\gdef\theauthnote{#1}}
  \newenvironment{frontmatter}{}{}
  \newenvironment{keywords}{\par\medskip\noindent\textbf{Keywords:} }{\par\medskip}
  \date{}
\fi
\usepackage{graphicx}
\usepackage{booktabs}
\usepackage{float}

% Generated by manuscript/make_numbers.py from experiment/runs; do not edit.

\newcommand{\EoneModels}{4}  % config.E1.models
\newcommand{\EoneReps}{10}  % config.E1.replicates
\newcommand{\EoneGenomes}{40}  % models x replicates
\newcommand{\EoneChroms}{10}  % config.E1.genome
\newcommand{\EoneChromMb}{50}  % config.E1.genome
\newcommand{\EoneGenomeMb}{500}  % config.E1.genome
\newcommand{\EoneRsCRmseMin}{\ensuremath{3.7\times10^{-6}}}  % E1_rs_vs_c rmse
\newcommand{\EoneRsCRmseMax}{\ensuremath{2.0\times10^{-4}}}  % E1_rs_vs_c rmse
\newcommand{\EoneRunawayC}{\ensuremath{1.1\times10^{6}}}  % E1_curves worst replicate, C psmc's first lambda
\newcommand{\EoneLoglikeFirstRelMax}{\ensuremath{7.0\times10^{-13}}}  % E1_rs_vs_c loglike_iteration0_rel_diff
\newcommand{\EoneFoldMin}{1,500}  % E1: min over replicates of (smaller error vs truth) / rmse rs vs C, floored to 2 digits
\newcommand{\EoneFoldMedian}{7,000}  % E1: median rs error vs truth / median rmse rs vs C (Fig 1e)
\newcommand{\EoneConstantTruthRs}{\ensuremath{0.048}}  % E1_table1 rmse_truth_rs_mean
\newcommand{\EoneBottleneckTruthRs}{\ensuremath{0.291}}  % E1_table1 rmse_truth_rs_mean
\newcommand{\EtwoRsCRmseMax}{\ensuremath{2.8\times10^{-4}}}  % E2_rs_vs_c
\newcommand{\EtwoHumanSequences}{22}  % E2_fit_summary NA12878 sequences
\newcommand{\EtwoTurtleSeqMax}{1,220}  % E2_fit_summary turtle sequences
\newcommand{\EtwoTurtles}{3}  % E2_fit_summary turtle genomes
\newcommand{\EtwoBootReps}{100}  % E2_bootstrap replicates
\newcommand{\EtwoBootThreads}{50}  % E2_bootstrap threads
\newcommand{\EtwoBandOverDeltaMin}{260}  % E2_bootstrap band_width_log10_median / E2_rs_vs_c rmse, smallest over the genomes, floored to 2 digits
\newcommand{\EtwoSplitFits}{6}  % E2_split_patterns split rows
\newcommand{\EtwoSplitGainMin}{\ensuremath{49}}  % E2_split_patterns loglike_gain_vs_default
\newcommand{\EtwoSplitGainMax}{\ensuremath{586}}  % E2_split_patterns loglike_gain_vs_default
\newcommand{\EthreeRepeats}{5}  % config.E3.repeats
\newcommand{\EthreeBigMb}{3,000}  % config.E3.genome_sizes_mb
\newcommand{\EthreeCWall}{4,456}  % E3_table2 c 1 thread
\newcommand{\EthreeCMin}{74.3}  % E3_table2 c 1 thread (min)
\newcommand{\EthreeRsOneWall}{317}  % E3_table2 rs02 1 thread
\newcommand{\EthreeRsOneMin}{5.28}  % E3_table2 rs02 1 thread (min)
\newcommand{\EthreeSpeedOne}{\ensuremath{14.1}}  % E3_table2 rs02 1 speedup
\newcommand{\EthreeSpeedEight}{\ensuremath{101}}  % E3_table2 rs02 8 speedup
\newcommand{\EthreeSpeedSixteen}{\ensuremath{198}}  % E3_table2 rs02 16 speedup
\newcommand{\EthreeSpeedThirty}{\ensuremath{386}}  % E3_table2 rs02 30 speedup
\newcommand{\EthreeRsWallThirty}{\ensuremath{11.5}}  % E3_table2 rs02 30 wall
\newcommand{\EthreeCMem}{1,062}  % E3_table2 c peak MB
\newcommand{\EthreeRsMemOne}{\ensuremath{36}}  % E3_table2 rs02 1 peak MB
\newcommand{\EthreeRsMemThirty}{\ensuremath{148}}  % E3_table2 rs02 30 peak MB
\newcommand{\EthreeSpeedOneMin}{\ensuremath{13.9}}  % E3_table2 rs02 1 thread speedup, every size
\newcommand{\EthreeSpeedOneMax}{\ensuremath{14.2}}  % E3_table2 rs02 1 thread speedup, every size
\newcommand{\EthreeChromWallMin}{\ensuremath{10.46}}  % E3_table2 rs02 one 100 Mb chromosome, any threads
\newcommand{\EthreeChromWallMax}{\ensuremath{10.97}}  % E3_table2 rs02 one 100 Mb chromosome, any threads
\newcommand{\EthreeEffMin}{\ensuremath{0.88}}  % E3_thread_scaling efficiency, 2+ threads
\newcommand{\EthreeEffMax}{\ensuremath{0.97}}  % E3_thread_scaling efficiency, 2+ threads
\newcommand{\EthreeNodes}{4}  % E3_runs hostname
\newcommand{\EthreeSetups}{19}  % E3_table2 rows, C psmc and psmc-rs
\newcommand{\EthreeSpreadMax}{4.4\%}  % E3_table2 (max-min)/median wall time, the 3000 Mb setups of Fig 2
\newcommand{\EthreeDetRsTotal}{75}  % E3_determinism rs02
\newcommand{\EthreeDetCTotal}{20}  % E3_determinism c
\newcommand{\EthreeCpu}{AMD EPYC 7702}  % E3_environment machine.cpu
\newcommand{\EfourDispThirtyFqTheta}{34.9\%}  % E4_routes dispersed30 fq2psmcfa -theta_rel (median)
\newcommand{\EfourDispThirtyAutoTheta}{1.8\%}  % E4_routes dispersed30 mask_auto -theta_rel (median)
\newcommand{\EfourDispTenFqTheta}{11.9\%}  % E4_routes dispersed10 fq2psmcfa -theta_rel (median)
\newcommand{\EfourDispTenAutoTheta}{0.7\%}  % E4_routes dispersed10 mask_auto -theta_rel (median)
\newcommand{\EfourDispThirtyMhsTheta}{31.9\%}  % E4_routes dispersed30 multihetsep -theta_rel (median)
\newcommand{\EfourDispThirtyNomaskTheta}{35.5\%}  % E4_routes dispersed30 nomask -theta_rel (median)
\newcommand{\EfourClusThirtyMhsTheta}{9.1\%}  % E4_routes clustered30 multihetsep -theta_rel (median)
\newcommand{\EfourClusThirtyNomaskTheta}{56.8\%}  % E4_routes clustered30 nomask -theta_rel (median)
\newcommand{\EfourReps}{5}  % config.E4.replicates
\newcommand{\EfourDispTenPartly}{48\%}  % E4_summary dispersed10 mask_auto
\newcommand{\EfourDispThirtyPartly}{92\%}  % E4_summary dispersed30 mask_auto
\newcommand{\EfourClusRunKb}{\ensuremath{20}}  % config.E4.masks mean uncallable run
\newcommand{\EfourDispRunBp}{\ensuremath{20}}  % config.E4.masks mean uncallable run
\newcommand{\EfourGenomeMb}{500}  % config.E4.genome
\newcommand{\EfiveFits}{50}  % E5_summary fits
\newcommand{\EfiveDatasets}{10}  % E5_summary datasets
\newcommand{\EfiveReps}{50}  % E5_summary replicates_per_dataset
\newcommand{\EfiveChroms}{6}  % config.E5.genome
\newcommand{\EfiveFiftySd}{\ensuremath{0.84}}  % E5_summary sd_ratio_median
\newcommand{\EfiveFiftyCov}{37\%}  % E5_summary coverage_mean
\newcommand{\EfiveFiftyCovMean}{83\%}  % E5_summary coverage_of_fit_mean_mean
\newcommand{\EfiveFiftyTrunks}{6}  % E5_summary trunks
\newcommand{\EfiveFiveSd}{\ensuremath{0.93}}  % E5_summary sd_ratio_median
\newcommand{\EfiveFiveCov}{43\%}  % E5_summary coverage_mean
\newcommand{\EfiveFiveCovMean}{92\%}  % E5_summary coverage_of_fit_mean_mean
\newcommand{\EfiveFiveTrunks}{60}  % E5_summary trunks
\newcommand{\EfiveBias}{\ensuremath{0.061}}  % E5_summary abs_bias_log10_median
\newcommand{\EfiveSdTrue}{\ensuremath{0.028}}  % E5_summary sd_true_log10_median
\newcommand{\EsixMhsRmseMax}{\ensuremath{6.7\times10^{-3}}}  % E6_S2_fits multihetsep rmse_log10_ne_vs_psmcfa_per_interval, max over E1's first replicates
\newcommand{\EsixDecodeCov}{95.6\%}  % E6_S3_summary coverage95_states
\newcommand{\EsevenGenomes}{2,504}  % E7_samples rows
\newcommand{\EsevenPops}{26}  % E7_samples populations
\newcommand{\EsevenDefault}{52}  % E7_samples default
\newcommand{\EsevenPerPopMin}{61}  % E7_samples per population
\newcommand{\EsevenPerPopMax}{113}  % E7_samples per population
\newcommand{\EsevenDepth}{\ensuremath{36.2}}  % E7_inputs raw_depth median
\newcommand{\EsevenMissingShare}{7.6\%}  % E7_route_shift missing share, median
\newcommand{\EsevenThetaExpected}{8.2\%}  % 1 / (1 - missing share) - 1
\newcommand{\EsevenThetaShift}{8.5\%}  % E7_route_shift theta_rel median
\newcommand{\EsevenRouteShiftMedian}{\ensuremath{0.044}}  % E7_route_shift mean_log10_ne_vcf_vs_psmcfa, median
\newcommand{\EsevenRouteShapeMedian}{\ensuremath{0.008}}  % E7_route_shift sqrt(rmse^2 - mean^2) of log10 Ne vcf vs psmcfa (the change of shape), median
\newcommand{\EsevenSplitSpanMedian}{\ensuremath{1.1}}  % E7_split log10(max / min) of the four split lambdas' Ne, median
\newcommand{\EsevenRsCRmseMin}{\ensuremath{5.7\times10^{-6}}}  % E7_rs_vs_c
\newcommand{\EsevenRsCRmseMax}{\ensuremath{2.1\times10^{-5}}}  % E7_rs_vs_c
\newcommand{\RealGenomes}{55}  % E2_rs_vs_c and E7_rs_vs_c samples (NA12878 in both)
\newcommand{\EsevenGiabPrecision}{\ensuremath{0.9992}}  % E7_giab precision
\newcommand{\EsevenGiabRecall}{\ensuremath{0.9966}}  % E7_giab recall
\newcommand{\EsevenGiabMb}{2,109}  % E7_giab eval_bases (Mb)
\newcommand{\EsevenRsCoreHours}{257}  % E7_throughput rs-psmcfa core_hours
\newcommand{\EsevenRsFitMin}{6.22}  % E7_throughput rs-psmcfa wall median (min)
\newcommand{\EsevenCFitMin}{88.0}  % E7_throughput c wall median (min)
\newcommand{\EsevenCCoreHoursAll}{3,700}  % E7_throughput c wall median x genomes (estimate)
\newcommand{\EsevenPairedSpeedMedian}{\ensuremath{13.8}}  % E7_fits: C psmc's wall time over psmc-rs's on the same genome, median over the comparison subset
\newcommand{\EsevenBootReps}{100}  % E7_bootstrap replicates
\newcommand{\EsevenBootThreads}{32}  % E7_bootstrap threads
\newcommand{\EsevenMu}{\ensuremath{1.25\times10^{-8}}}  % config.E7.scaling
\newcommand{\EsevenG}{29}  % config.E7.scaling
\newcommand{\EsevenBootFiles}{103}  % E2/bootstrap and E7/boot NA12878 files
\newcommand{\EsevenBootFilesIdentical}{103}  % E2/bootstrap and E7/boot NA12878 files byte-identical
\newcommand{\EsevenWindowLoKy}{\ensuremath{10.6}}  % E7_pop_bands all window_lo_years (ky)
\newcommand{\EsevenWindowHiMy}{\ensuremath{4.6}}  % E7_pop_bands all window_hi_years (My)
\newcommand{\EsevenNonAfrMinNeLo}{1,961}  % E7_pop_bands non-AFR superpopulation minima
\newcommand{\EsevenNonAfrMinNeHi}{2,902}  % E7_pop_bands non-AFR superpopulation minima
\newcommand{\EsevenNonAfrMinKyLo}{\ensuremath{43}}  % E7_pop_bands non-AFR minima (ky)
\newcommand{\EsevenNonAfrMinKyHi}{\ensuremath{52}}  % E7_pop_bands non-AFR minima (ky)
\newcommand{\EsevenAfrMinNe}{9,422}  % E7_pop_bands AFR minimum
\newcommand{\EsevenConvergeKy}{\ensuremath{247}}  % E7_pop_bands: superpopulation medians within 20% from here on
\newcommand{\EsevenPopCutKy}{\ensuremath{150}}  % split of recent and older times for the populations' spread
\newcommand{\EsevenPopSpreadRecentLo}{\ensuremath{0.24}}  % E7_pop_bands populations' max log10 spread, first window
\newcommand{\EsevenPopSpreadRecentHi}{\ensuremath{0.99}}  % E7_pop_bands populations' max log10 spread, first window
\newcommand{\EsevenPopSpreadMidLo}{\ensuremath{0.19}}  % E7_pop_bands populations' max log10 spread, after it to 150 ky
\newcommand{\EsevenPopSpreadMidHi}{\ensuremath{0.24}}  % E7_pop_bands populations' max log10 spread, after it to 150 ky
\newcommand{\EsevenPopSpreadOldMax}{\ensuremath{0.11}}  % E7_pop_bands populations' max log10 spread, >= 150 ky
\newcommand{\EsevenPopSpreadMidMedLo}{\ensuremath{0.015}}  % E7_pop_bands populations' log10 spread after the first window to 150 ky, median over times
\newcommand{\EsevenPopSpreadMidMedHi}{\ensuremath{0.084}}  % E7_pop_bands populations' log10 spread after the first window to 150 ky, median over times
\newcommand{\EsevenPopSpreadOldMed}{\ensuremath{0.015}}  % E7_pop_bands populations' log10 spread >= 150 ky, median over times, largest over the superpopulations
\newcommand{\EsevenIndivSpreadLo}{\ensuremath{0.035}}  % E7_pop_bands superpopulations but AMR: log10(ne_hi / ne_lo), 10% to 90% of genomes, median over the times after the first window
\newcommand{\EsevenIndivSpreadHi}{\ensuremath{0.040}}  % E7_pop_bands superpopulations but AMR: log10(ne_hi / ne_lo), 10% to 90% of genomes, median over the times after the first window
\newcommand{\EsevenRecords}{123,574,606}  % MANIFEST chromosomes records
\newcommand{\EsevenNonpass}{11,256,876}  % MANIFEST chromosomes nonpass
\newcommand{\EsevenPilotChrom}{21}  % experiment/e7_pilot/pilot_chr21.json
\newcommand{\EsevenLongestMb}{249}  % experiment/real_data NA12878 psmcfa (E7's input): longest sequence, bins x bin size (Mb)
\newcommand{\AgreeIntervals}{4,620}  % every interval in the window of every genome both tools fitted
\newcommand{\AgreeGenomes}{95}  % genomes both tools fitted (NA12878 once)
\newcommand{\AgreeMedianAbs}{\ensuremath{5.5\times10^{-6}}}  % every interval in the window of every genome both tools fitted: median |log10 Ne difference|
\newcommand{\AgreeMaxAbs}{\ensuremath{6.6\times10^{-4}}}  % every interval in the window of every genome both tools fitted: largest |log10 Ne difference|
\newcommand{\AgreeMaxPct}{0.15\%}  % every interval in the window of every genome both tools fitted: largest relative Ne difference
\newcommand{\EoneSoneGains}{1,920}  % Fig S1: gains drawn (genomes x programs x iterations)
\newcommand{\EoneSoneGainLastMin}{\ensuremath{0.43}}  % Fig S1: gains in the last iteration drawn
\newcommand{\EoneSoneGainLastMax}{\ensuremath{5.9}}  % Fig S1: gains in the last iteration drawn
\newcommand{\EoneSoneLastIter}{24}  % Fig S1: last iteration both programs report
\newcommand{\EoneSoneFirstAbsMax}{\ensuremath{4.9\times10^{-7}}}  % E1_traces: largest |LL_rs - LL_c| at the start
\newcommand{\EoneSoneRsHigher}{23}  % E1_traces: genomes whose psmc-rs log-likelihood is higher at the last common iteration
\newcommand{\EoneSoneDiffMax}{\ensuremath{0.065}}  % E1_traces: largest |LL_rs - LL_c| after any iteration
\newcommand{\EoneSoneRatioMin}{84}  % E1_traces: a model's smallest gain in iteration i over its largest |LL_rs - LL_c| after i, smallest over the models and iterations (Fig S1's panels), floored to 2 digits
\newcommand{\EoneSoneThetaMax}{\ensuremath{3.9\times10^{-4}}}  % E1_traces and E1_rs_vs_c theta_rel_diff: largest relative theta difference, any iteration
\newcommand{\EoneSoneRhoMax}{\ensuremath{3.6\times10^{-4}}}  % E1_traces and E1_rs_vs_c rho_rel_diff: largest relative rho difference, any iteration
\newcommand{\DiagOverCBuilds}{3}  % c_equivalence/summary.json default rmse_vs_c median / c_builds_rmse median, rounded: the median Delta in units of C psmc's own noise
\newcommand{\DiagMedianDefault}{\ensuremath{2.5\times10^{-5}}}  % c_equivalence/summary.json default rmse_vs_c median
\newcommand{\DiagMedianLastBin}{\ensuremath{6.2\times10^{-6}}}  % c_equivalence/summary.json lastbin rmse_vs_c median
\newcommand{\DiagMedianHookeJeeves}{\ensuremath{2.6\times10^{-5}}}  % c_equivalence/summary.json hj rmse_vs_c median
\newcommand{\DiagCBuildsMedian}{\ensuremath{8.5\times10^{-6}}}  % c_equivalence/summary.json c_builds_rmse median
\newcommand{\DiagLastBinCloser}{37}  % c_equivalence/summary.json lastbin closer_to_c_than_default
\newcommand{\DiagLastBinOther}{3}  % c_equivalence/summary.json lastbin genomes not closer
\newcommand{\DiagPrecisionMedian}{\ensuremath{1.0\times10^{-7}}}  % c_equivalence/summary.json prec rmse_vs_default median
\newcommand{\DiagPrecisionMax}{\ensuremath{2.5\times10^{-6}}}  % c_equivalence/summary.json prec rmse_vs_default max
\newcommand{\XpRsRmseMax}{\ensuremath{2.4\times10^{-6}}}  % crossplatform rs RMSE window max (E1)
\newcommand{\XpParamRelMax}{\ensuremath{2.1\times10^{-5}}}  % crossplatform rs param_max_rel max (E1): largest relative difference of any parameter
\newcommand{\XpRatioCMedian}{\ensuremath{124}}  % side_by_side ratio C cross-build / rs cross-platform, median
\newcommand{\EoneModelsW}{four}  % config.E1.models (in words)
\newcommand{\EtwoTurtlesW}{three}  % E2_fit_summary turtle genomes (in words)
\newcommand{\EtwoSplitFitsW}{six}  % E2_split_patterns split rows (in words)
\newcommand{\EthreeRepeatsW}{five}  % config.E3.repeats (in words)
\newcommand{\EthreeNodesW}{four}  % E3_runs hostname (in words)
\newcommand{\EfourRepsW}{five}  % config.E4.replicates (in words)
\newcommand{\EfiveChromsW}{six}  % config.E5.genome (in words)
\newcommand{\EfiveFiftyTrunksW}{six}  % E5_summary trunks (in words)
\newcommand{\DiagOverCBuildsW}{three}  % c_equivalence/summary.json default rmse_vs_c median / c_builds_rmse median, rounded: the median Delta in units of C psmc's own noise (in words)
\newcommand{\DiagLastBinOtherW}{three}  % c_equivalence/summary.json lastbin genomes not closer (in words)

\newcommand{\psmcrs}{psmc-rs}
\newcommand{\Ne}{\ensuremath{N_\mathrm{e}}}
\newcommand{\opt}[1]{\texttt{#1}}

\journal{gigabyte}
\papercat{Technical Release}
\runningauthor{Han}

\title{psmc-rs: fast PSMC inference that reproduces the reference implementation}
\author[1,\authfn{1}]{Yufan Han}
\affil[1]{School of Life Sciences, Fudan University, Shanghai, China}
\authnote{\authfn{1}larryivanhan@gmail.com}

\begin{document}

\begin{frontmatter}
\maketitle
\ifdefined\gigabytepreview\noindent{\small *Correspondence: larryivanhan@gmail.com}\par\fi

\begin{abstract}
The pairwise sequentially Markovian coalescent (PSMC) infers how the effective size of a population changed over time from the genome of one diploid individual. Its reference implementation, the C program psmc, runs on one thread and takes more than an hour per human genome. psmc-rs reimplements PSMC in Rust: it keeps C psmc's model and estimator, and fits the chromosomes of a genome in parallel. On \EoneGenomes{} simulated and \RealGenomes{} real genomes, its estimates of population size differed from C psmc's by at most \AgreeMaxPct{} over the time range each genome informs, far less than the method's own error. It was \EthreeSpeedOne{} times as fast as C psmc on one thread and \EthreeSpeedThirty{} times as fast on 30 threads, used a fraction of the memory, and gave the same results on any number of threads. Fitting all \EsevenGenomes{} unrelated high-coverage genomes of the 1000 Genomes Project took \EsevenRsCoreHours{} core-hours.

\textbf{Availability and implementation:} psmc-rs is written in Rust and released under the MIT license at \url{https://github.com/larryivan/psmc-rs}; version 0.2.0 is archived at \url{https://doi.org/10.5281/zenodo.23762893}. The inputs and fits of the \EsevenGenomes{} genomes are at \url{https://doi.org/10.5281/zenodo.23764975}.
\end{abstract}

\begin{keywords}
demographic inference; effective population size; PSMC; coalescent hidden Markov model; population genomics; reproducibility; software
\end{keywords}
\end{frontmatter}

\section{Statement of Need}
\label{sec:background}

The history of a population's size leaves a trace in a single genome. Along the two copies of a diploid genome, the time to their most recent common ancestor (TMRCA) changes at recombination events, and it is shorter where the population was small. The pairwise sequentially Markovian coalescent \citep[PSMC;][]{li2011} reads this trace: it treats the TMRCA as the hidden state of a hidden Markov model whose observations are short bins of sequence, each heterozygous or homozygous, and fits a piecewise-constant history of the effective population size (\Ne) by expectation--maximization (EM). Because it needs only one unphased genome, PSMC has been applied to modern and archaic humans \citep{li2011,prufer2014,mallick2016} and to a wide range of other species \citep{nadachowska2016,hilgers2025}, and it remains the baseline for newer coalescent hidden Markov models \citep{schiffels2014,terhorst2017,cousins2025,terhorst2025}.

PSMC's reference implementation, the C program \texttt{psmc} \citep{psmcC}, runs on one thread. Twenty-five EM iterations on a 3~Gb genome take more than an hour, a bootstrap multiplies that time by its number of replicates, and population resources such as the high-coverage 1000 Genomes data \citep{byrskabishop2022} now hold thousands of genomes. Other programs port PSMC to other languages and toolkits \citep{psmcr,psmcpython,smckit} or extend its model \citep{schiffels2020,liu2022,cousins2024,cousins2025}. To replace C psmc, however, a faster program must keep its estimator and reproduce its results, so that new curves remain comparable with the many PSMC curves already published. The native PSMC of smckit has been checked against C psmc on one chromosome \citep{smckit}; we are not aware of a program shown to reproduce C psmc's fits of whole genomes on the same inputs while running substantially faster than C psmc.

\psmcrs{} is that program. It fits the same model as C psmc, from the same inputs and with the same options, an order of magnitude faster on one thread and two orders of magnitude faster on many threads for genomes of many sequences, and its results do not depend on the number of threads. It is meant for anyone who fits PSMC: to fit one genome quickly, to bootstrap it, or to fit a cohort of thousands of genomes.

\section{Implementation}

\paragraph{Model} \psmcrs{} fits the PSMC model of \citet{li2011} with the conventions of C psmc 0.6.5-r74: the same discretization of time, the same transition and emission probabilities (its code ports C psmc's \texttt{psmc\_update\_hmm} line by line), the same parameter patterns and starting point, and EM on the same objective. It differs from C psmc in four ways. First, it processes the sequences of a genome (its chromosomes or scaffolds) in parallel, each on one thread. Second, its maximization step (M-step) runs L-BFGS \citep{liu1989} to convergence within bounds on the parameters, where C psmc runs a Hooke--Jeeves search \citep{hooke1961}. Third, its expected emission counts include the last bin of each sequence, as both programs' log-likelihoods do; C psmc leaves that bin out of its counts. Fourth, when the input records how many of a bin's bases are callable, that is, could be genotyped, it weights the bin by that number (below).

\paragraph{Expectation step} The forward--backward algorithm runs in double precision on each sequence. To bound memory, \psmcrs{} keeps the forward vector only every 128 bins and recomputes the vectors in between, to the same bits, during the backward pass; C psmc keeps the vectors of every bin. For models of up to 256 states the forward and backward passes take time linear in the number of states, from the rank-one structure of PSMC's transition matrix \citep{harris2014}, where C psmc's take time quadratic in it. At start-up \psmcrs{} chooses kernels compiled for x86-64 processors with AVX2 and FMA, NEON kernels on aarch64, or portable ones otherwise.

\paragraph{Determinism} The expectation step and the log-likelihood run in parallel over groups of whole sequences whose grouping depends only on the number of sequences, and their partial sums are merged in a fixed order, so results do not depend on the number of threads. The block bootstrap \citep{kunsch1989} follows C psmc's: it cuts each sequence into blocks of 500,000 bins, as C psmc's \texttt{splitfa} does, and each replicate draws blocks with replacement. Replicates run in parallel, each drawing from its own random stream seeded from the bootstrap seed and its index, so a seed gives the same replicates on any number of threads.

\paragraph{Partly callable bins} C psmc's input marks a 100~bp bin missing only when more than 90\% of its bases are uncallable, and fits a bin with fewer uncallable bases as if all its bases had been observed, which biases PSMC's estimates \citep{lipson2015,nadachowska2016}. When the input records the number $c$ of a bin's $s$ bases that are callable, as a VCF read with its callable mask does, \psmcrs{} gives a homozygous bin the probability $e^{-\theta\bar t_kc/s}$ instead of $e^{-\theta\bar t_k}$ in state $k$, the Poisson bin likelihood of recent coalescent hidden Markov models \citep{cousins2024,cousins2025} restricted to PSMC's binary observations. Inputs without callable counts are fitted with C psmc's model unchanged, and \opt{--bin-callable ignore} fits a VCF as C psmc's input would be fitted.

\paragraph{Inputs and outputs} \psmcrs{} reads C psmc's \texttt{psmcfa} files; VCF files, either all-sites consensus calls or variant calls with a mask of callable sites; and the multihetsep files of MSMC \citep{schiffels2014}. It writes the fitted model as JSON and as an HTML report and, on request, in C psmc's \texttt{.psmc} format, which C psmc's \texttt{psmc\_plot.pl} draws. It also converts inputs (\opt{psmc-rs convert}), checks them before fitting (\opt{psmc-rs qc}) and decodes the posterior distribution of the TMRCA in every bin (\opt{--tmrca-out}). With the settings of C psmc's documentation, this command fits a human genome for 25 EM iterations on eight threads:
\begin{verbatim}
psmc-rs NA12878.psmcfa.gz NA12878.json 25 --pattern "4+25*2+4+6" --t-max 15 \
    --tr-ratio 5 --mu 1.25e-8 --gen-years 29 --threads 8 --psmc-out NA12878.psmc
\end{verbatim}

\section{Validation}

\paragraph{Experiments} All results come from one pipeline, in the repository (\texttt{experiment/}), that simulates the genomes, runs both programs with matched settings, records how every output was made and draws every figure from its tables; \texttt{experiment/REPRODUCIBILITY.md} describes how to reproduce them. C psmc ran \texttt{psmc -N25 -t15 -r5 -p "4+25*2+4+6"}, the settings \citet{li2011} used for human genomes, and \psmcrs{} 0.2.0 the same settings with its own options, with 100~bp bins. We compared the two programs on \EoneGenomes{} simulated genomes of \EoneGenomeMb{}~Mb, \EoneReps{} under each of \EoneModelsW{} demographic models, simulated with msprime \citep{kelleher2016,baumdicker2022} (experiment E1); on \EtwoTurtlesW{} turtle genomes \citep{hilgers2025} and the human genome NA12878 (E2); and on \EsevenDefault{} human genomes of the 1000 Genomes Project, two per population, which we call the comparison subset (E7). Two fits are compared over the genome's informative window, the times between the 1\% and 99\% quantiles of its pairwise coalescence time, by $\Delta$, the root-mean-square error (RMSE) of one fit's $\log_{10}\Ne$ against the other's, interval by interval and weighted by the logarithmic length of each interval in the window. The fits ran on a computing cluster of AMD EPYC 7702 nodes.

\paragraph{Agreement with C psmc} On every genome, $\Delta$ was at most \EtwoRsCRmseMax{} (Fig.~\ref{fig:agree}, Table~\ref{tab:agreement}). Across all \AgreeIntervals{} intervals in the windows of these \AgreeGenomes{} genomes, the two programs' \Ne{} differed by a median of \AgreeMedianAbs{} in $\log_{10}\Ne$ and by at most \AgreeMaxPct{}. This difference is far smaller than the error of the method itself: on the simulated genomes, $\Delta$ was at least \EoneFoldMin{} times smaller than either program's error relative to the true history, and the two programs' errors were equal to three decimals. It is about \DiagOverCBuildsW{} times the difference between C psmc's own fits of the simulated genomes on two platforms (median \DiagCBuildsMedian{}), and far smaller than the sampling uncertainty of a fit: on the turtles and NA12878, the median width of \psmcrs{}'s bootstrap 95\% band was at least \EtwoBandOverDeltaMin{} times $\Delta$. A diagnostic traced the remaining difference mostly to the last bin of each sequence: leaving it out of \psmcrs{}'s expected emission counts, as C psmc does, lowered the median $\Delta$ on the simulated genomes from \DiagMedianDefault{} to \DiagMedianLastBin{}, whereas replacing L-BFGS with C psmc's search barely changed it (\DiagMedianHookeJeeves{}). In both programs the log-likelihood never decreased from one EM iteration to the next, and the two programs' log-likelihoods differed by at most \EoneSoneDiffMax{} units after any iteration (Additional file~1, Fig.~S1).

\paragraph{Speed and memory} We timed both programs on simulated constant-size genomes of 1, 3, 10 and 30 chromosomes of 100~Mb, each run alone on an exclusive node, with \EthreeRepeatsW{} repeats (E3; Fig.~\ref{fig:perf}, Table~\ref{tab:timing}). On one thread, 25 EM iterations on the \EthreeBigMb{}~Mb genome took \psmcrs{} \EthreeRsOneMin{} minutes and C psmc \EthreeCMin{} minutes: \psmcrs{} was \EthreeSpeedOne{} times as fast, and \EthreeSpeedOneMin{} to \EthreeSpeedOneMax{} times as fast at every size. With the sequences in parallel, it was \EthreeSpeedEight{}, \EthreeSpeedSixteen{} and \EthreeSpeedThirty{} times as fast as C psmc on 8, 16 and 30 threads, with a parallel efficiency of \EthreeEffMin{} to \EthreeEffMax{}. Because each sequence runs on one thread, a genome's longest sequence bounds the speed-up: one 100~Mb chromosome took \EthreeChromWallMin{} to \EthreeChromWallMax{}~s on 1, 8 and 30 threads alike, and a human genome is bounded in the same way by its chromosome~1 (\EsevenLongestMb{}~Mb). \psmcrs{} needed \EthreeRsMemOne{}~MB of memory on one thread and \EthreeRsMemThirty{}~MB on 30, against C psmc's \EthreeCMem{}~MB.

\paragraph{Thread and platform independence} For each genome size, all fitted models, from every thread count, repeat and node, were byte-identical (\EthreeDetRsTotal{} fits in all). NA12878's bootstrap, run with the same seed on \EtwoBootThreads{} and on \EsevenBootThreads{} threads, wrote \EsevenBootFilesIdentical{} of its \EsevenBootFiles{} files byte for byte the same. Refitted on a laptop of another processor architecture, operating system and math library, the simulated genomes' parameters differed from the cluster's fits by at most a relative \XpParamRelMax{}.

\paragraph{Partly callable bins and the bootstrap} On simulated genomes with 10\% or 30\% of positions uncallable in short runs (E4), taking every called bin as full underestimated the scaled mutation rate $\theta$, and with it every size and time, by \EfourDispTenFqTheta{} or \EfourDispThirtyFqTheta{}, and the weighted fit by only \EfourDispTenAutoTheta{} or \EfourDispThirtyAutoTheta{} (medians over five genomes; Additional file~1, Fig.~S2). In a calibration study of one bottleneck model (E5), the bootstrap's 95\% bands held the mean of independent fits at, on average, \EfiveFiveCovMean{} of the times in the informative window with blocks of 5~Mb and \EfiveFiftyCovMean{} with blocks of 50~Mb, but the true history at only \EfiveFiveCov{} and \EfiveFiftyCov{} (Additional file~1, Fig.~S3): the bands measure the sampling variation of an estimate, not PSMC's bias.

\paragraph{Application to the 1000 Genomes Project} We fitted \psmcrs{} to all \EsevenGenomes{} unrelated individuals of the high-coverage 1000 Genomes data \citep{byrskabishop2022}, \EsevenPops{} populations sequenced to a median depth of \EsevenDepth{}$\times$, from one \texttt{psmcfa} file per genome made with the strict accessibility mask \citep{kgp2015} (E7; Fig.~\ref{fig:kg}). Fitting one genome per thread on a shared node took \psmcrs{} a median of \EsevenRsFitMin{} minutes per genome and \EsevenRsCoreHours{} core-hours for the cohort; C psmc, run the same way on the comparison subset, took a median of \EsevenCFitMin{} minutes per genome, about \EsevenCCoreHoursAll{} core-hours for the cohort. On the comparison subset, $\Delta$ between the two programs was \EsevenRsCRmseMin{} to \EsevenRsCRmseMax{}. The histories have the features of published single-genome PSMC analyses of human genomes \citep{li2011,prufer2014,mallick2016}: the median curves of the four superpopulations outside Africa fell to minima of \EsevenNonAfrMinNeLo{} to \EsevenNonAfrMinNeHi{} between \EsevenNonAfrMinKyLo{} and \EsevenNonAfrMinKyHi{} thousand years ago, whereas the African median fell only to \EsevenAfrMinNe{}. Within each superpopulation, the population medians differed most over the span of PSMC's first free $\lambda$, a single parameter for the four most recent intervals that one genome constrains poorly, and stayed close after it. The \texttt{psmcfa} inputs, the fits and the provenance of every fit are released as a resource.

\paragraph{Limitations} We tested one pattern of parameters and 25 EM iterations, which do not reach convergence (Additional file~1, Fig.~S1), on genomes of up to \EtwoTurtleSeqMax{} sequences; among the four genomes of E2, the difference between the two programs grew with the number of sequences, so fits of highly fragmented assemblies are worth comparing with C psmc's. A VCF read with its callable mask is fitted with weighted bins by default, which C psmc's model does not do. \psmcrs{} keeps PSMC's limits: one genome says little about the last few thousand generations, where the first free $\lambda$ can run away (\psmcrs{} bounds it and warns), and ancestral population structure can produce changes in a single-genome curve that are not changes in size \citep{mazet2016,cousins2025}. Because \psmcrs{} parallelizes across sequences, a genome assembled into a few long sequences gains less from many threads than one assembled into many. The same machinery, a checkpointed expectation step run in parallel across sequences with results that do not depend on the number of threads, could also serve MSMC2 \citep{schiffels2020}, whose composite likelihood sums PSMC-like hidden Markov models over pairs of haplotypes.

\begin{figure*}[tbp]
\centering
\includegraphics[width=\textwidth]{art1_agreement.pdf}
\caption{\textbf{\psmcrs{} reproduces C psmc.} (a--d) Histories fitted by both programs to the same genome, C psmc's (wide, pale orange) under \psmcrs{}'s (thin, blue), over the genome's informative window. Below each, \psmcrs{}'s $\log_{10}\Ne$ minus C psmc's, interval by interval, in units of $10^{-4}$; $\Delta$, the RMSE of this difference over the window. (a, b) The first simulated genome of the bottleneck (a) and zigzag (b) models, with the true history (dashed); $\mu=2.5\times10^{-8}$ per bp per generation, $g=25$ years. (c, d) The human genome NA12878 ($\mu=1.25\times10^{-8}$, $g=29$ years) and the turtle \textit{Pelusios castaneus} ($\mu=4.6\times10^{-9}$, $g=15$ years), the genome of the largest $\Delta$, with, in pale blue where it is wider than the curves, the pointwise 95\% band of \psmcrs{}'s bootstrap (\EtwoBootReps{} replicates in blocks of 50~Mb; for the turtle, its \EtwoTurtleSeqMax{} whole scaffolds). Panel d and its strip have their own scales. The turtle's recent peak is an artefact of the pattern's single free $\lambda$ for the four most recent intervals \citep{hilgers2025}. (e) $\Delta$ for every genome both programs fitted: the \EsevenDefault{} human genomes of the comparison subset (NA12878 among them), the \EtwoTurtlesW{} turtles and the \EoneGenomes{} simulated genomes; rings, the genomes of a--d. The row ``C psmc on two platforms'' gives, for comparison, $\Delta$ between C psmc's own fits of the simulated genomes on two platforms. Below, each program's RMSE against the true history of the simulated genomes. Bars, medians of rows of at least ten genomes; the arrow joins the simulated genomes' median $\Delta$ to the median of \psmcrs{}'s error on them: $\Delta$ is \EoneFoldMedian{} times smaller.}
\label{fig:agree}
\end{figure*}

\begin{table*}[tbp]
\centering\footnotesize
\caption{\textbf{Agreement between \psmcrs{} and C psmc on every genome both programs fitted, with the same input and settings.} NA12878 is in both the E2 and the E7 rows. $\Delta$, the RMSE (range over genomes), and largest absolute difference $|d_k|$ of $\log_{10}\Ne$, interval by interval over the informative window; largest relative differences of the scaled mutation rate $\theta$, the scaled recombination rate $\rho$ and the maximum coalescence time $t_{\max}$, of the log-likelihood at the last iteration both programs report, and of the log-likelihood of the first expectation step, at the shared starting point, where the two agree at about the rounding of the six decimals C psmc prints.}
\label{tab:agreement}
\setlength{\tabcolsep}{3.5pt}
\begin{tabular}{lrrrrrr}
\toprule
& & \multicolumn{2}{c}{$\log_{10}N_\mathrm{e}$ in the window} & \multicolumn{3}{c}{largest relative difference} \\
\cmidrule(lr){3-4}\cmidrule(lr){5-7}
Genomes & $n$ & $\Delta$ & largest $|d_k|$ & $\theta,\rho,t_{\max}$ & log-likelihood & first E-step \\
\midrule
Simulated (E1) & 40 & \ensuremath{3.7\times10^{-6}} to \ensuremath{2.0\times10^{-4}} & \ensuremath{6.6\times10^{-4}} & \ensuremath{5.5\times10^{-4}} & \ensuremath{1.3\times10^{-8}} & \ensuremath{7.0\times10^{-13}} \\
Turtles and NA12878 (E2) & 4 & \ensuremath{9.7\times10^{-6}} to \ensuremath{2.8\times10^{-4}} & \ensuremath{4.5\times10^{-4}} & \ensuremath{1.0\times10^{-3}} & \ensuremath{2.3\times10^{-7}} & \ensuremath{6.9\times10^{-14}} \\
1000 Genomes (E7) & 52 & \ensuremath{5.7\times10^{-6}} to \ensuremath{2.1\times10^{-5}} & \ensuremath{5.8\times10^{-5}} & \ensuremath{5.5\times10^{-5}} & \ensuremath{2.8\times10^{-9}} & \ensuremath{1.1\times10^{-13}} \\
\bottomrule
\end{tabular}

\end{table*}

\begin{figure*}[tbp]
\centering
\includegraphics[width=\textwidth]{art2_performance.pdf}
\caption{\textbf{\psmcrs{} is faster than C psmc and needs less memory.} (a, b) A simulated constant-size genome of 30 chromosomes of 100~Mb (E3), fitted for 25 EM iterations, each run alone on an exclusive \EthreeCpu{} node; points are medians of \EthreeRepeatsW{} repeats on \EthreeNodesW{} nodes, whose wall times spread by at most \EthreeSpreadMax{} of their median. C psmc runs on one thread only: its square, at one thread, is drawn flat across the thread counts. (a) Wall time; arrows, the speed-up over C psmc on 1 and on 30 threads; grey line, perfect scaling of \psmcrs{}'s one-thread time; dotted line, \psmcrs{}'s time on one 100~Mb chromosome alone. (b) Peak memory (MB of $2^{20}$ bytes). (c) The 1000 Genomes cohort (E7), each genome fitted on one thread of a node shared with other fits: wall time per genome of C psmc on the \EsevenDefault{} genomes of the comparison subset and of \psmcrs{} on all \EsevenGenomes{}; bars, medians; arrow, the ratio of the medians (on the \EsevenDefault{} genomes both programs fitted, the median per-genome ratio was \EsevenPairedSpeedMedian{}). These times are not comparable with those of a.}
\label{fig:perf}
\end{figure*}

\begin{table*}[tbp]
\centering\small
\caption{\textbf{Run time and memory of C psmc and \psmcrs{} (E3).} Median wall time of 25 EM iterations over \EthreeRepeatsW{} repeats on exclusive \EthreeCpu{} nodes, for simulated constant-size genomes of 100~Mb chromosomes; speed-up over C psmc (a ratio of unrounded medians) and peak memory (median, MB of $2^{20}$ bytes) on \EthreeBigMb{}~Mb. -- marks a setup not timed.}
\label{tab:timing}
\begin{tabular}{lrrrrrrr}
\toprule
& & \multicolumn{4}{c}{Wall time, median (s)} & \multicolumn{2}{c}{3,000 Mb} \\
\cmidrule(lr){3-6}\cmidrule(lr){7-8}
Tool & Threads & 100 Mb & 300 Mb & 1,000 Mb & 3,000 Mb & Speed-up & Memory (MB) \\
\midrule
C psmc & 1 & 152.5 & 449.4 & 1,478.3 & 4,455.9 & 1 & 1,062 \\
psmc-rs & 1 & 11.0 & 31.8 & 103.9 & 317.0 & 14.1 & 36 \\
psmc-rs & 2 & -- & -- & -- & 163.3 & 27.3 & 39 \\
psmc-rs & 4 & -- & -- & -- & 87.3 & 51.0 & 47 \\
psmc-rs & 8 & 10.5 & 11.0 & 21.5 & 43.9 & 101 & 59 \\
psmc-rs & 16 & -- & -- & -- & 22.5 & 198 & 88 \\
psmc-rs & 30 & 10.5 & 11.0 & 11.5 & 11.5 & 386 & 148 \\
\bottomrule
\end{tabular}

\end{table*}

\begin{figure*}[tbp]
\centering
\includegraphics[width=\textwidth]{fig7_1000_genomes.pdf}
\caption{\textbf{Population size histories of all \EsevenGenomes{} unrelated individuals in the high-coverage 1000 Genomes data, each genome fitted by \psmcrs{} (E7).} (a) The median over each superpopulation's genomes. (b--f) One superpopulation each: its genomes (faint lines, which merge into a soft band where they agree), the median over each of its populations' genomes (\EsevenPerPopMin{} to \EsevenPerPopMax{} genomes per population; 1000 Genomes population codes; *, admixed) and, in grey, the other four superpopulations' medians from a. Medians are pointwise, at each time of a common grid, and span the cohort's informative window, from \EsevenWindowLoKy{} thousand to \EsevenWindowHiMy{} million years ago. Pale: PSMC's first free $\lambda$, a single parameter for the four most recent intervals, which one genome constrains poorly; each median is pale until nine in ten of its genomes have passed it. The fits are of the \texttt{psmcfa} inputs, which take every called bin as full; weighting bins by their callable bases raised the curves of the comparison subset almost uniformly, by a median of \EsevenRouteShiftMedian{} in $\log_{10}\Ne$. $\mu=\EsevenMu$ per bp per generation, $g=\EsevenG$ years.}
\label{fig:kg}
\end{figure*}

\section{Availability of source code and requirements}
\begin{itemize}
\item Project name: psmc-rs
\item Project home page: \url{https://github.com/larryivan/psmc-rs}
\item Operating system(s): Linux and macOS (Windows through WSL)
\item Programming language: Rust
\item Other requirements: none for the release binaries; Rust 1.88 or later to build from source
\item License: MIT
\item Archived version: psmc-rs 0.2.0, \url{https://doi.org/10.5281/zenodo.23762893}; also on crates.io
\end{itemize}

\section{Availability of Test Data}
The test inputs and expected outputs of \psmcrs{}'s test suite, including C psmc's own outputs on the same inputs, are in the repository (\texttt{tests/data}). The experiment pipeline (\texttt{experiment/}) regenerates the simulated genomes from fixed seeds, and its committed results hold every fit, bootstrap and table of this article; the turtle \texttt{psmcfa} files, from \citet{hilgers2025}, and NA12878's are in \texttt{experiment/real\_data}. The inputs of the 1000 Genomes analysis, the \texttt{psmcfa} files of the \EsevenGenomes{} genomes, the comparison subset's VCFs and the masks, are on Zenodo (\url{https://doi.org/10.5281/zenodo.23764975}). The 1000 Genomes high-coverage VCFs are available from the International Genome Sample Resource (\url{https://www.internationalgenome.org}).

\section{Declarations}

\subsection{List of abbreviations}
EM: expectation--maximization; MB: megabyte ($2^{20}$ bytes); PSMC: pairwise sequentially Markovian coalescent; RMSE: root-mean-square error; TMRCA: time to the most recent common ancestor; VCF: variant call format.

\subsection{Ethical Approval}
Not applicable. The human genomes analysed are openly released data of the 1000 Genomes Project.

\subsection{Competing Interests}
The author declares no competing interests.

\subsection{Funding}
This work received no specific funding.

\subsection{Author's Contributions}
Yufan Han: conceptualization, methodology, software, validation, formal analysis, investigation, data curation, visualization, and writing (original draft, review and editing).

\subsection{Use of AI tools}
The text of this manuscript was polished with the assistance of AI. The author takes full responsibility for the results and the content of the article.

\section{Acknowledgements}
We thank Dr Ke Wang and the members of her laboratory at Fudan University for their support. We thank the 1000 Genomes Project, the New York Genome Center, Genome in a Bottle and \citet{hilgers2025} for making their data available.

\bibliography{references}

\end{document}
